% =====================================================================
% Discussion generale et conclusion : texte suivi, sans intitules de
% section, conformement au plan valide. Le mecanisme reprend celui de
% l'environnement conclusion de la classe memoireuqam (chapitre non
% numerote, entree de table des matieres), avec le titre retenu.
% =====================================================================
\chapter*{DISCUSSION GÉNÉRALE ET CONCLUSION}
\addtocontents{toc}{\protect\vspace{1.5ex}}
\addcontentsline{toc}{chapter}{DISCUSSION GÉNÉRALE ET CONCLUSION}

% Ce chapitre n'etant pas numerote, le compteur de chapitre n'avance pas et
% ses tableaux heriteraient du numero du chapitre 6. On leur donne un prefixe
% propre, comme l'annexe utilise A. La numerotation par chapitre est retablie
% en fin de fichier, pour que l'annexe retrouve A.1 et suivants.
\setcounter{secnumdepth}{0}   % intitules non numerotes : \section* est
% Les intitules de ce chapitre ne portent pas de numero. Sans lui, la classe
% UQAM les compose en romain de taille normale et ils se confondent avec le
% corps du texte. Ils passent donc en italique, ici seulement : les sections
% numerotees des chapitres 1 a 6 gardent leur presentation, leur numero
% suffisant a les identifier. La definition d'origine est retablie plus bas,
% avec le compteur.
\makeatletter
\renewcommand\section{\@startsection{section}{1}{\z@}%
  {12bp\@plus.2ex\@minus.2ex}{0.01bp\@plus.2ex}%
  {\normalfont\normalsize\itshape\baselineskip=\simpleinter}}
\makeatother
                              % inoperant dans memoireuqam.cls
\setcounter{table}{0}
\renewcommand{\thetable}{D.\arabic{table}}

Ce mémoire s'est déployé en six temps. Le Chapitre~\ref{chapitre:contexte} a établi l'enjeu
économique et zootechnique de la boiterie, son étiologie multifactorielle, les limites de la
notation locomotrice conventionnelle et les variables comportementales que restituent les
capteurs portés. Le Chapitre~\ref{chapitre:etatart} a passé en revue les approches automatisées
de détection, les méthodes non supervisées de référence et la détection temporelle de
changements persistants, avant de montrer que l'hétérogénéité des protocoles interdit de
comparer directement les performances publiées. Le Chapitre~\ref{chapitre:systeme} a défini le
système proposé~: la période de référence individuelle, la branche HYPO, la branche
INSTABILITÉ, leur fusion hiérarchique, les comparateurs et le protocole d'évaluation, avec les
métriques et l'organisation logicielle qui les rendent reproductibles. Le
Chapitre~\ref{chapitre:donnees} a décrit le corpus IceQube 2023, la préparation des données et
le cadre expérimental. Le Chapitre~\ref{chapitre:interface} a présenté l'interface de
priorisation, ses trois niveaux de priorité et les limites de son usage en ferme. Le
Chapitre~\ref{chapitre:resultats} a enfin rapporté la performance attribuable du module HYPO,
sa comparaison aux détecteurs d'anomalies ponctuels, ses analyses de sensibilité, le test de
stress à aire d'enveloppe appariée, l'extension bidirectionnelle et la confrontation à la
cohorte IceTag-SLS.

\section{Contributions}
Au terme de ce travail, les contributions peuvent être résumées ainsi. Le mémoire fournit une
évaluation attribuable de HYPO avec intervalles de confiance, une ablation et ses tailles
d'effet, une analyse OFAT, une courbe reliant détection et charge, une tentative de
calibration individuelle rapportée malgré son résultat négatif, une extension bidirectionnelle
assortie de contrôles de confusion, une évaluation du temps d'exécution, une analyse SLS
synchronisée, une interface de priorisation et des artefacts reproductibles. Sa contribution
méthodologique est double~: proposer un détecteur temporel interprétable, et montrer comment
évaluer ce type d'outil sans confondre résultats techniques, hypothèses exploratoires et
interprétations cliniques. Le code source, les paramètres figés, les douze profils d'injection
et les scripts d'évaluation sont versionnés, avec un manifeste d'intégrité SHA-256 des sorties
et de l'empreinte du corpus source (Annexe~\ref{annexe:validation})~; les données capteurs brutes ne peuvent être diffusées pour
des raisons de confidentialité. Les artefacts et leurs empreintes publiés permettent d'auditer
les résultats, et la démarche complète, des injections jusqu'aux figures du manuscrit, reste
reproductible par les personnes autorisées à accéder aux données.

Le module HYPO constitue en définitive le résultat technique principal~: il démontre sa
capacité à transformer des baisses graduelles en épisodes attribuables, surpasse les variantes
IF/LOF évaluées pour la localisation et conserve une logique interprétable par vache. Sa
performance de localisation et sa charge de vérification définissent désormais des objectifs
d'amélioration mesurables.
L'extension bidirectionnelle organise une hypothèse complémentaire et quantifie ses compromis.
Cette séparation entre preuve principale et exploration rend l'ensemble cohérent et fournit
une base reproductible pour développer et évaluer un futur usage de priorisation en ferme.

\section{Ce que le système établit}
L'objectif de ce mémoire était de concevoir une chaîne reproductible qui transforme des
mesures indirectes du capteur IceQube en priorités de vérification locomotrice. Son apport central est
de montrer qu'une alerte comportementale peut être construite autour d'une trajectoire
individuelle plutôt qu'autour de points anormaux isolés. La branche HYPO combine direction du
changement, cohérence entre signaux et persistance temporelle. Cette combinaison produit
des épisodes plus cohérents avec les dégradations graduelles que les comparateurs IF
et LOF. Après soustraction de l'exécution propre, HYPO couvre 29 des 33 événements
graduels et atteint IoU20 dans 13 cas. Ce résultat établit la capacité du système à
transformer une dégradation contrôlée en priorité de vérification localisée.

Le résultat établi se lit à deux niveaux complémentaires. La couverture de 87,9\,\% montre que HYPO
réagit dans la plupart des fenêtres injectées, tandis que le taux de nouveaux départs de
57,6\,\% distingue les épisodes réellement ouverts de ceux qui prolongent une alerte déjà
active. Le délai médian de 20,75 heures situe le caractère \textit{précoce} relativement aux
dégradations de 36 à 60 heures. Ces indicateurs précisent l'usage opérationnel du système~:
prioriser une vérification lorsque plusieurs familles de signaux se dégradent de façon
persistante, avec une traçabilité suffisante pour expliquer pourquoi une vache remonte dans le
classement. Le Tableau~\ref{tab:portee_conclusions} associe chaque conclusion à son niveau de
preuve.


La comparaison de HYPO aux autres détecteurs apporte un résultat méthodologique important~: la
quantité de points anormaux n'est pas le critère décisif. IF ponctuel produit
7,658 départs par vache-jour sans atteindre IoU20 sur un seul événement, et les variantes IF
et LOF suivies de persistance réduisent la charge sans localiser davantage les fenêtres. HYPO
atteint, pour sa part, un IoU attribuable moyen de 0,137 et un taux de 29,5\,\% à IoU20 sur
les quatre profils. La direction du changement et l'accumulation temporelle sont donc mieux soutenues
que la simple accumulation d'anomalies ponctuelles. La contribution propre de la cohérence
entre familles demande toutefois une comparaison supplémentaire, car le comparateur
pédométrique partage la même machinerie temporelle.

\begin{table}[H]
  \centering
  \caption{Portée des conclusions principales.}
  \label{tab:portee_conclusions}
  \small
  \begin{tabular}{p{0.22\textwidth}p{0.34\textwidth}p{0.35\textwidth}}
    \toprule
    Niveau & Conclusion permise & Conclusion non permise \\
    \midrule
    Démontré techniquement
      & HYPO réagit à des baisses graduelles après la période de référence et localise mieux ces scénarios
        que les comparateurs IF/LOF.
      & HYPO identifie à lui seul la cause de la dégradation. \\
    Soutenu mais limité
      & Le mécanisme temporel reste robuste à cinq formes moins favorables à aire
        d'enveloppe appariée.
      & La cohérence multivariée apporte un gain démontré sur un comparateur pédométrique
        doté de la même machinerie temporelle. \\
    Exploratoire
      & INSTABILITÉ peut produire une sortie de surveillance distincte et la fusion peut
        représenter une séquence synthétique.
      & L'instabilité précède habituellement une boiterie réelle. \\
    À établir dans un protocole clinique
      & Validation longitudinale avec examens locomoteurs et causes documentées.
      & Sensibilité, spécificité, valeur prédictive ou bénéfice clinique. \\
    \bottomrule
  \end{tabular}
\end{table}

Le test de stress à aire d'enveloppe appariée apporte cette comparaison sur 5 formes moins
favorables, 3 positions et 11 vaches. Sur les 165 événements positifs, HYPO atteint en moyenne
97,0\,\% de couverture attribuable, 71,5\,\% de nouveaux départs et un taux de 30,9\,\% à
IoU20. HYPO conserve un avantage exact sur IF, IF avec persistance et LOF avec persistance
($p=0{,}000488$ pour chaque comparaison). Ce résultat positif montre que l'avantage de HYPO
sur les détecteurs ponctuels ne dépend pas uniquement de la forme graduelle initiale. En
revanche, sa couverture ne dépasse pas celle du comparateur pédométrique temporel
(différence $-0{,}012$, $p=0{,}844$), et le contrôle pas-seuls échoue au critère principal de
spécificité: 78,8\,\% de recouvrement attribuable et 30,3\,\% de nouveaux départs. Ce contrôle
reste toutefois mal localisé (IoU moyen de 0,038 et aucun événement à IoU20). La campagne
soutient donc la robustesse temporelle de HYPO aux formes moins favorables, mais ne démontre
pas un gain indépendant de spécificité multivariée.

L'analyse OFAT montre ensuite que ces conclusions ne reposent pas sur \textbf{une valeur unique} de
chaque seuil. Elle identifie toutefois deux choix structurants~: l'agrégation et le nombre de
familles concordantes. Exiger quatre familles réduit fortement la couverture et le fond, alors
que n'en exiger que deux augmente la charge~; une \textbf{agrégation courte} réagit plus vite mais
localise moins bien, tandis qu'une \textbf{agrégation longue} perd des événements. Les valeurs de
référence constituent donc un compromis documenté, non un optimum. Modifier les seuils après
lecture de cette analyse aurait transformé une étude de robustesse en sélection sur les
données d'évaluation~; aucun réglage a posteriori n'a été effectué.

Deux analyses complémentaires permettent aussi d'écarter des choix moins favorables. La
calibration par la dispersion du score propre à
chaque période de référence ne réduit ni le fond moyen ni sa variabilité, et diminue IoU20 de
39,4 à 33,3\,\%~: elle est rejetée. La courbe reliant la détection à la charge montre par
ailleurs qu'un seuil de 0,14 paraîtrait favorable sur cette campagne, mais l'adopter après
observation des mêmes événements introduirait un biais de sélection. Ce résultat sert à
formuler une hypothèse de développement ultérieur, non à améliorer rétroactivement la
performance annoncée.

L'extension bidirectionnelle répondait à une question secondaire~: peut-on représenter une
instabilité courte sans l'assimiler automatiquement à une alerte de vérification~? Sur le plan
logiciel, l'architecture y répond grâce à deux persistances indépendantes et à une sortie de
surveillance distincte. Les contrôles brefs et l'activité coordonnée de type œstrus sont
rejetés dans la campagne synthétique, la séquence instabilité vers hypoactivité est récupérée
dans 90,9\,\% des cas, et 66,7\,\% des instabilités produisent bien une sortie de
surveillance. La fusion hiérarchique augmente la couverture des scénarios de vérification de
61,4 à 72,7\,\% par rapport à HYPO seul dans la campagne étendue, qui compte 132 scénarios. En contrepartie, elle
n'améliore pas IoU20, elle porte le fond de 0,446 à 0,571 notification par vache-jour et elle
déclenche une alerte sur une part notable des facteurs confondants. Son apport démontré réside donc dans la séparation
des sorties et dans le gain de couverture, plutôt que dans une supériorité globale sur HYPO.
Les profils d'instabilité et leur chronologie restent des hypothèses de test~: la branche
structure ces hypothèses dans une interface de surveillance, tandis que HYPO demeure la
méthode principale.

Le rejet des 33 contrôles brefs vérifie qu'un pic capteur, deux heures d'exercice synthétique
ou une manipulation d'une heure ne suffisent pas à déclencher le détecteur dans les conditions
testées. Cette réussite dépend toutefois de la représentation choisie et ne garantit pas le
rejet de tous les événements réels de même nature. La confusion avec l'hypoactivité non
locomotrice constitue quant à elle la limite centrale du travail~: neuf scénarios sur onze,
soit 81,8\,\%, sont alertés, car leur forme est volontairement identique à une dégradation
locomotrice sur les signaux disponibles, et 40,9\,\% de l'ensemble des facteurs confondants
déclenchent une notification. Aucun réglage temporel ne peut séparer entièrement deux causes
qui produisent les mêmes observations. La température, l'ingestion, la production laitière,
les traitements, le vêlage ou un examen clinique sont nécessaires pour améliorer la
spécificité causale. Les capteurs reconnaissent une forme comportementale sans en identifier
l'origine, ce qui impose une vérification terrain.

\section{Charge de vérification et pistes d'amélioration}
La charge que ce dispositif représente mérite d'être regardée en face. Le fond de 0,402
notification par vache-jour pour HYPO et de 0,571 pour la fusion hiérarchique demeure élevé
pour un grand troupeau. Ces valeurs ne sont pas des taux de faux positifs, puisqu'une donnée
propre peut contenir une dégradation réelle non annotée, mais elles indiquent une fatigue
d'alerte potentielle. Un déploiement devrait regrouper les alertes par épisode, limiter les
répétitions, prioriser les trajectoires persistantes et enregistrer le résultat de la
vérification humaine.

Deux leviers supplémentaires, absents du système actuel, méritent d'être signalés car ils
portent sur une dimension différente. Le premier est une normalisation à l'échelle du
troupeau~: le détecteur compare aujourd'hui chaque vache à sa seule période de référence, ce
qui le rend insensible à un décalage affectant simultanément tout le troupeau. Le second est
un filtre des journées collectives~: lorsqu'une proportion élevée d'animaux alerte le même
jour, la cause la plus probable relève de la conduite, du logement ou du climat plutôt que
d'une atteinte individuelle. La carte thermique présentée par la Figure~\ref{fig:streamlit_heatmap} rend
déjà ces journées visibles à l'écran, mais aucun mécanisme ne les écarte du décompte.

Ces deux leviers se distinguent du filtre d'activité coordonnée déjà présent dans la branche
INSTABILITÉ, qui compare l'excès de pas à celui de l'indice de mouvement chez une même vache
et demeure donc intra-individuel. Ils agiraient sur une dimension indépendante des trois niveaux de
priorité~: là où la priorité indique quel mécanisme a déclenché, une normalisation par le
troupeau et un filtre collectif permettraient de graduer la confiance accordée à chaque
alerte, et donc de réduire la charge de revue sans toucher au détecteur lui-même. Cette piste
n'est pas évaluée dans ce mémoire, et son effet sur la couverture attribuable reste à mesurer~: écarter
une journée collective peut aussi écarter une dégradation individuelle qui y coïncide.

Du côté du coût de calcul, le temps médian de 10,64 secondes pour
28 vaches montre que le prototype est techniquement exécutable sur le corpus étudié~;
IF représente la majeure partie de ce coût, tandis que le calcul
HYPO, INSTABILITÉ et fusion ne représente que 6,4\,\% du temps instrumenté. Ce résultat ne
préjuge pas de la latence en ferme, du stockage continu ou de la concurrence entre
utilisateurs.

\section{Confrontation aux scores locomoteurs et à la littérature}
L'analyse IceTag-SLS apporte un résultat encourageant~: le nombre de notifications dans les
sept jours précédant le score est plus élevé chez les trois vaches SLS~$\geq2$, avec une AUC
exploratoire de 0,924 et un test exact bilatéral global à $p=0{,}030$, tandis que le score
maximal et la surveillance d'instabilité séparent moins bien les groupes. L'AUC reste entre
0,909 et 0,955 lorsqu'une vache positive est retirée, et un test max-stat unilatéral
exploratoire demeure favorable parmi les cinq variantes de notifications ($p=0{,}0467$). En revanche, il devient
non concluant sur les vingt combinaisons variante-métrique ($p=0{,}0742$), ce qui borne la
portée inférentielle sans effacer le signal descriptif. Cette concordance
reste aussi limitée par l'effectif et par le fait que les trois vaches SLS~$\geq2$ appartiennent
au groupe \textit{Exercise}. Dans ce seul groupe, l'AUC est de 0,778 et le test exact
bilatéral est non concluant ($p=0{,}40$); aucune AUC interne n'est estimable dans le groupe
sans exercice, faute de cas SLS~$\geq2$. Le résultat appuie donc la faisabilité de l'analyse
synchronisée et un signal descriptif à approfondir, pas une validation clinique ni un
réglage des seuils.

Replacés dans la littérature, ces résultats s'accordent avec la signature comportementale la
mieux établie de la boiterie, à savoir une réduction de l'activité locomotrice et une
augmentation du temps couché \citep{ito2010lying,blackie2011lying,chapinal2010measurement}~;
la branche HYPO cible précisément cette direction. Les revues soulignent toutefois que les
variables utiles et leurs amplitudes dépendent du capteur, du protocole et du contexte
\citep{oleary2020accelerometers,alsaaod2019automatic}. La référence propre à chaque vache est
donc un choix méthodologique prudent pour absorber une partie de l'hétérogénéité observée, et
elle améliore la comparabilité des trajectoires individuelles. Le
Tableau~\ref{tab:revue_ciblee} du Chapitre~\ref{chapitre:etatart} récapitule les travaux
auxquels ces choix se rattachent.

Ce choix d'une référence propre à chaque vache, et à chaque créneau horaire, répond à une
structure de variance connue~: chez la vache laitière, l'activité varie fortement entre
individus (âge, rang de lactation, stade, tempérament) et selon le rythme circadien lié à la
traite et à l'alimentation. Un seuil de population risquerait donc d'être gouverné par
l'identité de la vache et par l'heure, plutôt que par son changement propre. En normalisant
chaque vache par sa médiane horaire de référence, HYPO réduit l'influence de cette variance
de nuisance interindividuelle et horaire, et rend plus comparable, d'une vache à l'autre,
l'ampleur d'une dégradation relative.

\section{Limites du capteur et comparaison aux systèmes existants}
Les limites du capteur doivent être rappelées avec la même franchise. L'agitation, les reports
de poids et les changements posturaux peuvent être observés lors d'un inconfort, mais les
agrégats IceQube ne mesurent ni l'appui relatif des membres ni la démarche. Une hausse de
mouvement ou de transitions reste en outre compatible avec l'œstrus, l'exercice, une
manipulation ou un artefact, ce qui explique que la branche INSTABILITÉ filtre l'activité
coordonnée et demeure exploratoire. Enfin, les performances rapportées varient fortement d'une
étude à l'autre \citep{oleary2020accelerometers,qiao2021review}, ce qui rappelle qu'aucune
signature unique n'est pleinement fiable sur ces capteurs et conforte le cadrage en alerte non
diagnostique.

Une comparaison avec les systèmes commerciaux appelle la même retenue. Plusieurs solutions
fondées sur des podomètres et des accéléromètres d'activité sont déployées en élevage, mais
leur logique interne et leurs protocoles d'évaluation ne sont pas toujours décrits avec le
niveau de détail nécessaire à une comparaison indépendante
\citep{rutten2013invited,stangaferro2016use}. Le présent travail se distingue par un code
auditable, des paramètres figés, un protocole explicite et un manifeste d'intégrité. Il ne
revendique aucune performance supérieure à ces systèmes, faute d'étalon et de protocole
communs.

Pour l'éleveur, à l'échelle d'un troupeau, le système se positionne comme un outil de
priorisation et non d'alarme. La charge de fond mesurée reste trop élevée pour une
vérification systématique~; un déploiement réaliste afficherait plutôt chaque jour les
quelques vaches les plus anormales, à partir du classement de l'interface. Le bénéfice attendu
est d'orienter plus tôt l'attention humaine vers des vaches dont le comportement dérive, le
coût étant le temps de vérification et le risque de fatigue d'alerte. Cette forme de tri
correspond mieux à la charge observée qu'une alarme binaire envoyée pour chaque notification.
Le Tableau~\ref{tab:forces_limites} synthétise les forces et les limites de l'ensemble.

\begin{table}[H]
  \centering
  \caption{Forces et limites du travail.}
  \label{tab:forces_limites}
  \small
  \begin{tabular}{p{0.43\textwidth}p{0.48\textwidth}}
    \toprule
    Forces & Limites \\
    \midrule
    Séparation temporelle stricte et exécution propre de référence.
      & Absence de labels cliniques dans le corpus principal. \\
    Profils graduels, multivariés et physiquement contraints.
      & Amplitudes partiellement hypothétiques. \\
    Ablation sur les mêmes événements et agrégation par vache.
      & Onze vaches seulement dans les campagnes. \\
    OFAT de 20 variantes pour HYPO.
      & Aucune interaction entre paramètres ni jeu de développement indépendant. \\
    Contrôles brefs et facteurs confondants explicites.
      & Hypoactivité locomotrice et non locomotrice non identifiables avec ces seuls signaux. \\
    Évaluation du temps d'exécution de la chaîne complète et artefacts reproductibles.
      & Temps mesuré sur une seule machine et un seul corpus. \\
    Analyse SLS fondée sur des mesures antérieures au score.
      & Trois SLS~$\geq2$, tous dans le groupe \textit{Exercise}. \\
    \bottomrule
  \end{tabular}
\end{table}

\section{Menaces à la validité}
Ces limites gagnent à être examinées selon la typologie classique des menaces à la validité en
recherche empirique \citep{shadish2002experimental,wohlin2012experimentation}. Du côté du
construit, les injections représentent des changements de comportement et non des cas de
boiterie~; les capteurs ne donnent aucune mesure directe de la douleur, de la lésion ou de
l'asymétrie de démarche, et le seuil IoU20 est une convention technique dépourvue de signification clinique.

Ce protocole d'injection soulève une objection légitime~: la campagne principale reproduit
une baisse graduelle, persistante et concordante, c'est-à-dire précisément la signature que
HYPO recherche. Trois éléments limitent toutefois le risque de circularité.
D'abord, HYPO n'est pas évalué seul, mais confronté à IF et à LOF sur exactement
les mêmes événements injectés~; comme ces comparateurs échouent à localiser les épisodes avec
le critère IoU20 alors que HYPO obtient une localisation supérieure, l'avantage mesuré ne
découle pas de la simple présence de l'événement, mais de la manière dont il est capté.
Ensuite, le contrôle bref, une variation isolée d'une seule famille, n'appartient pas à la
cible de HYPO et est correctement rejeté, ce qui montre que le détecteur ne réagit pas à toute
perturbation injectée. L'amplitude du profil modéré reprend par ailleurs un ordre de grandeur
rapporté dans la littérature pour des changements comportementaux associés à des épisodes
locomoteurs ou à un parage thérapeutique \citep{paudyal2021lying}. Enfin, le test de stress
réexécute la comparaison sur un départ abrupt, des familles désynchronisées, une récupération
asymétrique, du bruit et un bloc manquant. L'aire numérique sous l'enveloppe est appariée
entre les formes et le comparateur pédométrique est inclus; cet appariement ne garantit pas
une amplitude physique ou comportementale identique. Le test confirme que l'avantage sur IF
et LOF ne dépend pas de la seule forme graduelle. Il révèle aussi une limite importante: le
contrôle pas-seuls atteint 78,8\,\% de recouvrement attribuable, même s'il reste très mal
localisé. L'apport propre de la cohérence multivariée n'est donc pas isolé. Les formes
injectées demeurent des hypothèses d'évaluation technique et non des trajectoires cliniques
observées.

Sur le plan de la validité interne, la séparation chronologique et la soustraction de
l'exécution propre réduisent la fuite et l'attribution erronée. Plusieurs scénarios
proviennent toutefois de la même vache, de sorte que l'unité indépendante reste la vache~; la
graine unique améliore la traçabilité, mais ne mesure pas toute la variabilité possible du
placement des événements. La période de référence est en outre définie automatiquement par les
premiers 60\,\% des intervalles admissibles, sans adjudication clinique de leur normalité. Les
médianes par créneau horaire et la comparaison au jour précédent réduisent l'influence des
valeurs isolées et du rythme journalier, mais une maladie présente dès le début, une dérive
capteur, une activité exceptionnellement faible ou un changement de conduite peuvent encore
biaiser les niveaux attendus et la charge d'alerte. Un déploiement devrait donc vérifier la
stabilité et la couverture de cette période, permettre sa confirmation par l'utilisateur et
prévoir sa réinitialisation après un changement majeur de contexte. Sur le plan de la validité
externe, le corpus principal provient d'une seule ferme et d'une seule période~: les différences
de logement, de sol, de saison, de conduite et de capteur peuvent modifier les distributions,
et la cohorte SLS, issue du même site, ne constitue pas une validation inter-fermes.

La validité statistique appelle enfin une lecture nuancée. Les proportions synthétiques sont
descriptives, les tests d'ablation utilisent onze paires au niveau vache, ce qui respecte la
structure mais limite la puissance, et les analyses OFAT et INSTABILITÉ examinent plusieurs
variantes sans jeu de test indépendant. Le détecteur ayant été conçu à partir du même corpus,
la campagne principale constitue une évaluation technique interne, et non un test indépendant
sur des vaches ou un site restés entièrement étrangers à la conception. La configuration a
néanmoins été gelée avant cette campagne finale, la période surveillée est temporellement
disjointe de la période de référence et aucune valeur observée dans l'OFAT, la calibration
individuelle ou la courbe détection-charge n'est utilisée pour remplacer la configuration de
référence. Ces précautions limitent la sélection rétrospective et la fuite temporelle, sans
supprimer l'optimisme possible lié aux décisions de conception prises sur le même corpus. Une
estimation pleinement externe demande donc un troupeau indépendant, avec un jeu de
développement distinct du jeu d'évaluation. Les AUC et valeurs de $p$ de l'analyse SLS restent
par ailleurs fragiles en raison du petit effectif, de la multiplicité et de la confusion
avec le traitement. Le test exact interne au groupe \textit{Exercise}, plutôt que le seul
test global, borne l'interprétation clinique.

\section{Éthique, gouvernance et validation clinique}
De ces limites découle un usage prévu étroit et un ensemble d'obligations éthiques. L'usage
défendable est un tri des vaches à observer. Une notification signifie que plusieurs
dimensions du comportement ont changé de manière persistante par rapport à la période de
référence de la vache~; elle ne signifie pas que la vache soit boiteuse. Une vérification doit
inclure l'observation locomotrice, l'examen du pied si nécessaire et une consultation
vétérinaire selon le contexte. L'absence d'alerte ne garantit pas l'absence de douleur, et le
système ne doit automatiser ni traitement ni réforme.

La gouvernance des données mérite d'être explicitée en vue d'un déploiement. Les droits de
propriété et d'usage ne peuvent pas être présumés~: ils doivent être définis entre
l'exploitation, le fournisseur technologique et les chercheurs, avec des règles explicites de
consentement et de confidentialité \citep{wolfert2017big,carbonell2016ethics}. Dans ce
travail, l'extraction CowAlert demeure locale, est exclue des artefacts diffusables et ne sert
qu'au contrôle de provenance~; seules des sorties agrégées et leurs empreintes d'intégrité
sont conservées dans la chaîne reproductible. Un système déployé devrait documenter la
finalité, les droits d'accès, la sécurité, la durée de conservation et l'audit des données, et
garantir qu'une alerte n'entraîne aucune décision automatique de traitement ou de réforme.

Le chemin vers une revendication clinique est dès lors clairement balisé. Une validation
clinique devrait réunir des scores locomoteurs datés, plusieurs évaluateurs, des examens des
pieds ou diagnostics vétérinaires, les causes non locomotrices documentées, les métadonnées
de contexte et les données capteurs des jours précédents. Le
développement et le test devraient être séparés par vache, par période et, idéalement, par
ferme, ce qui permettrait de régler les seuils sur le jeu de développement puis de les geler
avant le test, sans optimisme. Des contrôles réels doivent inclure œstrus, exercice,
manipulation, vêlage, chaleur et maladies non locomotrices. Au-delà du protocole, des signaux
de démarche ou d'asymétrie seraient plus proches du construit clinique que la seule quantité
de mouvement, et la charge de vérification devra être évaluée avec les utilisateurs afin de
définir un compromis acceptable en ferme.

\section{Perspectives et travaux futurs}
Les résultats obtenus dessinent un programme de travail dont les priorités sont ordonnées par
ce que les données actuelles ne permettent pas d'établir.

La première est un corpus indépendant. Une estimation pleinement externe demande un troupeau
resté étranger à la conception, un jeu de développement distinct du jeu d'évaluation et des
scores locomoteurs relevés sur l'ensemble des animaux plutôt que sur une sous-cohorte. Ce
cadre seul autoriserait à mesurer une sensibilité et une spécificité cliniques, termes que ce
mémoire s'est interdit d'employer faute d'étalon vétérinaire.

La deuxième porte sur la charge de vérification. Le plafond quotidien d'alertes, laissé ouvert
au Chapitre~\ref{chapitre:resultats}, demande à être fixé sur des données de terrain et son
coût en épisodes manqués à être mesuré. La normalisation par le troupeau et le filtre collectif
évoqués plus haut relèvent du même chantier~: réduire l'effort de revue sans modifier le
détecteur lui-même.

La troisième concerne la localisation. Un IoU attribuable moyen de 0,137 laisse une marge
considérable, et l'analyse OFAT désigne l'agrégation et le nombre de familles concordantes
comme les deux leviers structurants. Les explorer suppose toutefois un jeu de développement
séparé, sans quoi tout gain observé resterait une sélection rétrospective.

Le coût de calcul offre par ailleurs une marge immédiate~: le comparateur IF consomme
82,8\,\% du temps d'exécution alors que les deux branches et leur fusion n'en demandent que
6,4\,\%. Le retirer d'un déploiement, où il ne sert plus de point de comparaison, ramènerait
la chaîne complète sous deux secondes.

La quatrième, enfin, est la validation clinique longitudinale que le
Tableau~\ref{tab:portee_conclusions} classe parmi ce qui reste à établir~: examens locomoteurs
répétés, causes documentées et suivi des décisions prises à partir des alertes. Elle seule
dirait si la priorisation proposée change quelque chose au repérage de la boiterie en ferme.

% Retablir la numerotation par chapitre : l'annexe qui suit doit repartir en A.1.
\renewcommand{\thetable}{{\thechapter.\arabic{table}}}

\setcounter{secnumdepth}{5}

% Retablir la presentation d'origine des titres de section : l'annexe qui suit
% numerote les siens et n'a donc pas besoin de l'italique pose plus haut.
\makeatletter
\renewcommand\section{\@startsection{section}{1}{\z@}%
  {12bp\@plus.2ex\@minus.2ex}{0.01bp\@plus.2ex}%
  {\normalfont\normalsize\baselineskip=\simpleinter}}
\makeatother
